\documentclass[11pt]{article}
\usepackage[a4paper,margin=23mm]{geometry}
\usepackage{amsmath,amssymb,xcolor,hyperref,fancyhdr}
\usepackage{fontspec}
\setmainfont{Times New Roman}
\setmonofont{Times New Roman}
\usepackage[italic]{mathastext}
\setlength{\emergencystretch}{3em}
\setlength{\parskip}{7pt}\setlength{\parindent}{0pt}
\pagestyle{fancy}\fancyhf{}\setlength{\headheight}{14pt}
\lhead{FlowMLLab}\rhead{Reading and experiment guide}\cfoot{\thepage}
\begin{document}
\noindent\textbf{Week 7.2 | Reading the evidence}\par\medskip\hrule\medskip
\section*{A sensor-space check is not a field-space uncertainty check}
The reduced state contains POD coefficients. The observation matrix maps those
coefficients to sensor readings. The innovation is the measured reading minus
its predicted reading; the Kalman gain determines how much that innovation
changes the state. A near-unit normalized innovation RMS tests uncertainty at
observed sensors. It does not certify coverage of the full field, where missing
POD modes and model discrepancy remain. Thus innovation RMS near 1.08 and field
coverage near 55\% need not contradict one another.

If rank 8, 32 sensors and covariance scale 1000 are the largest searched values
and also the selected values, the search has reached its boundary. Extend the
search on development/validation cases before freezing a new protocol. Do not
extend it in response to test-case scores.
\paragraph{Check your understanding.}
Identify the input information, the reference, the metric denominator and the
evaluation split in one worked notebook comparison. Predict which conclusion
would change if any one of those four changed, then check your prediction.
\end{document}
